\documentclass[12pt,a4paper]{article}

\usepackage[french]{babel}
\usepackage[utf8]{inputenc}
\usepackage[T1]{fontenc}
\usepackage{fouriernc}
\usepackage{geometry}
\usepackage{array}
\usepackage{booktabs}
\usepackage{tabularx}
\usepackage[table]{xcolor}
\usepackage{amsmath}
\usepackage{fancyhdr}
\usepackage{tikz}
\usepackage{pgfplots}
\usepackage[most]{tcolorbox}
\usepackage{enumitem}
\pgfplotsset{compat=1.18}

\geometry{top=2.2cm, bottom=2cm, left=2cm, right=2cm, headheight=15pt}

%------------------------------------------------------------
% Couleurs (calquées sur le modèle moteurs_electriques.pdf)
%------------------------------------------------------------
\definecolor{bluedark}{RGB}{50, 100, 160}
\definecolor{bluelight}{RGB}{210, 228, 248}
\definecolor{bluemid}{RGB}{170, 200, 235}
\definecolor{grisfond}{RGB}{245, 245, 245}

%------------------------------------------------------------
% En-tête / pied de page
%------------------------------------------------------------
\pagestyle{fancy}
\fancyhf{}
\lhead{\small Physique-Chimie -- Terminale}
\rhead{\small Propagation du son}
\cfoot{\thepage}
\renewcommand{\headrulewidth}{0.5pt}

%------------------------------------------------------------
% Boîtes tcolorbox
%------------------------------------------------------------
% Boîte de section (en-tête de partie)
\tcbset{
  stylepartie/.style={
    colback=bluedark,
    colframe=bluedark,
    coltext=white,
    fonttitle=\bfseries,
    boxrule=0pt,
    arc=3pt,
    left=6pt, right=6pt, top=4pt, bottom=4pt,
    before skip=10pt, after skip=6pt,
  },
  styledoc/.style={
    colback=bluelight,
    colframe=bluemid,
    boxrule=0.8pt,
    arc=3pt,
    left=8pt, right=8pt, top=5pt, bottom=5pt,
    before skip=6pt, after skip=4pt,
    fonttitle=\bfseries\small,
    attach boxed title to top left={yshift=-2mm, xshift=4mm},
    boxed title style={colback=bluemid, colframe=bluemid, arc=2pt},
  },
  styleinfo/.style={
    colback=grisfond,
    colframe=bluemid,
    boxrule=0.8pt,
    arc=3pt,
    left=8pt, right=8pt, top=5pt, bottom=5pt,
    before skip=6pt, after skip=6pt,
  },
}

%------------------------------------------------------------
% Commandes utiles
%------------------------------------------------------------
% Cadre pour les calculs (coins arrondis, style cohérent)
\newcommand{\cadrecalcul}[1][4cm]{%
  \vspace{5mm}%
  \begin{tcolorbox}[
    colback=white,
    colframe=bluemid,
    boxrule=0.8pt,
    arc=4pt,
    left=4pt, right=4pt,
    top=0pt, bottom=0pt,
    height=#1,
    after skip=2mm,
  ]
  \end{tcolorbox}%
}

% Lignes de réponse pointillées
\newcommand{\reponse}[1]{%
  \vspace{5mm}%
  \foreach \i in {1,...,#1}{%
    \par\noindent\makebox[\linewidth]{\dotfill}\vspace{2mm}%
  }%
}

% Question numérotée style ▶
\newcounter{numquestion}
\newcommand{\question}[1]{%
  \stepcounter{numquestion}%
  \vspace{6pt}%
  \noindent{\color{bluedark}$\blacktriangleright$~\textbf{Question~\thenumquestion.}}~#1%
  \vspace{2pt}%
}

%------------------------------------------------------------
\begin{document}
%------------------------------------------------------------

%============================================================
% TITRE
%============================================================
\begin{tcolorbox}[
  colback=white, colframe=bluedark, boxrule=1.5pt, arc=4pt,
  before skip=0pt, after skip=8pt,
  left=10pt, right=10pt, top=8pt, bottom=8pt,
]
  \begin{center}
    {\color{bluedark}\Large\bfseries Propagation du son}\\[4pt]
    {\large Activité documentaire}\\[6pt]
    \hrule height 0.4pt \vspace{6pt}
    \textbf{Durée :} 45 min \hfill \textbf{Niveau :} Terminale 
  \end{center}
\end{tcolorbox}

\vspace{2pt}
\noindent\rule{\linewidth}{0.5pt}
\vspace{4pt}

\noindent\textbf{Nom :} \dotfill \hspace{1cm} \textbf{Prénom :} \dotfill \hspace{1cm} \textbf{Classe :} \dotfill

\vspace{2pt}
\noindent\rule{\linewidth}{0.5pt}
\vspace{6pt}

%------------------------------------------------------------
% Compétences
%------------------------------------------------------------
\begin{tcolorbox}[styleinfo]
  \textbf{\color{bluedark}Compétences travaillées (référentiel BO)}
  \begin{itemize}[noitemsep, leftmargin=1.2em]
    \item Extraire des informations de documents et les exploiter.
    \item Utiliser la relation $v = d / \Delta t$ pour résoudre des problèmes.
    \item Relier les propriétés du son à des applications industrielles (contrôle non destructif).
  \end{itemize}
\end{tcolorbox}

%------------------------------------------------------------
% Contexte
%------------------------------------------------------------
\begin{tcolorbox}[styleinfo]
  \textbf{\color{bluedark}Contexte}\\[4pt]
  Dans un atelier de production mécanique, les techniciens TRPM utilisent régulièrement
  des appareils de \textbf{contrôle non destructif (CND)} pour vérifier l'intégrité des pièces
  sans les endommager. L'une des techniques les plus employées est le \textbf{contrôle par
  ultrasons}, qui repose sur les propriétés de propagation du son dans les matériaux.
  Dans cette activité, vous allez comprendre comment fonctionne cette technique à partir
  de l'étude de la propagation du son.
\end{tcolorbox}

\vspace{4pt}

%============================================================
% PARTIE A
%============================================================
\begin{tcolorbox}[stylepartie]
  Partie A -- Nature et propriétés du son \hfill {\normalfont\small (15 min)}
\end{tcolorbox}

%------------------------------------------------------------
% Document 1
%------------------------------------------------------------
\begin{tcolorbox}[styledoc, title={Document 1 -- La propagation du son}]
  Le son est une \textbf{onde mécanique} : il correspond à la propagation d'une perturbation
  de pression dans un milieu matériel (solide, liquide ou gaz). \textbf{Le son ne se propage
  pas dans le vide.}

  \medskip

  Le son se propage sous forme d'une \textbf{onde longitudinale} : les molécules du milieu
  vibrent dans la même direction que la direction de propagation, créant alternativement
  des zones de \textbf{compression} et de \textbf{dépression}.

  \medskip

  La vitesse de propagation du son s'appelle la \textbf{célérité}, notée $v$ (en m/s). Elle
  dépend du milieu traversé : elle est plus faible dans les gaz, plus élevée dans les
  liquides, et encore plus élevée dans les solides. Elle augmente aussi avec la
  \textbf{température} : à 0°C, $v \approx 331$~m/s dans l'air ; à 20°C, $v \approx 340$~m/s.

  \medskip

  \textbf{Relation fondamentale :} si un son parcourt une distance $d$ (en m) pendant
  une durée $\Delta t$ (en s) à la célérité $v$ (en m/s) :
  \[
    \boxed{v = \frac{d}{\Delta t}}
  \]
  \vspace{5mm}
\end{tcolorbox}

\newpage


\question{Le son peut-il se propager dans le vide ? Justifier en citant le document~1.}
\reponse{2}

\question{Comparer la célérité du son dans les gaz, les liquides et les solides. Dans quel milieu le son se propage-t-il le plus vite ?}
\reponse{2}

\question{La célérité du son dans l'air à 20°C est $v = 340$~m/s. Convertir cette valeur en km/h. \textit{(Rappel : $1$~m/s $= 3{,}6$~km/h)}}
\reponse{2}

\vspace{4pt}

%============================================================
% PARTIE B
%============================================================
\begin{tcolorbox}[stylepartie]
  Partie B -- Célérité du son dans les matériaux et application industrielle \hfill {\normalfont\small (15 min)}
\end{tcolorbox}

%------------------------------------------------------------
% Document 2
%------------------------------------------------------------
\begin{tcolorbox}[styledoc, title={Document 2 -- Célérité du son dans différents milieux et contrôle non destructif}]

  \textbf{Tableau des célérités du son :}

  \vspace{6pt}
  \begin{center}
  \footnotesize
  \renewcommand{\arraystretch}{1.3}
  \begin{tabular}{|l|c|c|c|c|c|c|c|}
    \hline
    \rowcolor{bluemid}
    \textbf{Milieu}          & \textbf{Air (20°C)} & \textbf{Eau (20°C)} & \textbf{Aluminium} & \textbf{Acier} & \textbf{Cuivre} & \textbf{Fonte} & \textbf{Caoutchouc} \\
    \hline
    \textbf{État}            & Gaz & Liquide & Solide & Solide & Solide & Solide & Solide \\
    \hline
    \textbf{Célérité (m/s)}  & 340 & 1\,480  & 6\,300 & 5\,900 & 4\,700 & 3\,500 & 150 \\
    \hline
  \end{tabular}
  \end{center}

  \vspace{8pt}

  \textbf{Le contrôle non destructif (CND) par ultrasons :}

  \medskip

  Un \textbf{capteur piézoélectrique} émet une courte impulsion ultrasonore dans la pièce
  à contrôler. L'onde se propage dans le matériau et se réfléchit sur les interfaces
  (fond de la pièce, fissure, inclusion\ldots). Le capteur reçoit les \textbf{échos} et mesure
  le temps de parcours $\Delta t$.

  En connaissant la célérité $v$ du son dans le matériau, on calcule la \textbf{distance
  parcourue}. Comme l'onde fait un trajet aller-retour :
  \[
    d = \frac{v \times \Delta t}{2}
  \]

  \begin{center}
  \begin{tikzpicture}[scale=0.88]
    \fill[gray!25] (0,0) rectangle (9,2);
    \draw[thick] (0,0) rectangle (9,2);
    % Défaut
    \fill[white] (6,0.7) ellipse (0.45 and 0.28);
    \draw[thick] (6,0.7) ellipse (0.45 and 0.28);
    \node[below] at (6,0.42) {\small défaut};
    % Capteur
    \fill[bluedark!80] (2.2,2) rectangle (3,2.6);
    \draw[thick] (2.2,2) rectangle (3,2.6);
    \node[above] at (2.6,2.62) {\small capteur};
    % Flèche fond (aller)
    \draw[->, thick, blue!80!black] (2.5,2) -- (2.5,0.05);
    \node[right, blue!80!black] at (2.55,1) {\small émission};
    % Flèche fond (retour)
    \draw[->, thick, red!70!black, dashed] (2.7,0.05) -- (2.7,2);
    \node[left, red!70!black] at (2.65,1.3) {\small écho fond};
    % Flèche défaut (aller)
    \draw[->, thick, blue!50!black] (2.85,2) -- (5.58,0.9);
    % Flèche défaut (retour)
    \draw[->, thick, red!50!black, dashed] (6.42,0.9) -- (2.95,2);
    \node[above right, red!50!black] at (5.5,1.5) {\small écho défaut};
    % Légende épaisseur
    \draw[<->, thick] (8.5,0) -- (8.5,2);
    \node[right] at (8.55,1) {\small $e$};
  \end{tikzpicture}
  \end{center}
\end{tcolorbox}

\newpage

\question{Relever dans le tableau la célérité du son dans l'acier et dans l'aluminium. Laquelle est la plus élevée ?}
\reponse{2}

\question{Pourquoi parle-t-on de contrôle \textit{non destructif} ? Quel est l'intérêt pour un technicien TRPM ?}
\reponse{3}

\question{Dans la formule $d = \dfrac{v \times \Delta t}{2}$, expliquer pourquoi on divise par 2.}
\reponse{2}

\vspace{4pt}

%============================================================
% PARTIE C
%============================================================
\begin{tcolorbox}[stylepartie]
  Partie C -- Application : contrôle d'une pièce en acier \hfill {\normalfont\small (15 min)}
\end{tcolorbox}

%------------------------------------------------------------
% Document 3
%------------------------------------------------------------
\begin{tcolorbox}[styledoc, title={Document 3 -- Oscillogramme d'un contrôle ultrasonore}]

  Un technicien TRPM effectue le contrôle d'une pièce en \textbf{acier} d'épaisseur
  nominale $e = 60$~mm. L'appareil enregistre les échos reçus en fonction du temps.

  \vspace{6pt}
  \begin{center}
  \begin{tikzpicture}[scale=0.7]
    \begin{axis}[
      width=13.5cm, height=5cm,
      xlabel={Temps $\Delta t$ ($\mu$s)},
      ylabel={Amplitude (u.a.)},
      xmin=0, xmax=25,
      ymin=-0.1, ymax=1.5,
      xtick={0,2,4,6,8,10,12,14,16,18,20,22,24},
      ytick={0,0.5,1},
      grid=major,
      grid style={dashed, gray!40},
      axis lines=left,
      tick label style={font=\small},
      label style={font=\small},
    ]
      % Impulsion initiale
      \addplot[blue!80!black, very thick] coordinates {(0,0)(0,1.3)(0.5,1.3)(0.5,0)};
      \node[above, blue!80!black, font=\small] at (axis cs:0.25,1.32) {émission};
      % Écho du défaut (écho A)
      \addplot[red!70!black, thick] coordinates {(8.1,0)(8.1,0.72)(8.6,0.72)(8.6,0)};
      \node[above, red!70!black, font=\small] at (axis cs:8.35,0.74) {écho A};
      % Écho du fond (écho B)
      \addplot[orange!80!black, thick] coordinates {(20.3,0)(20.3,0.95)(20.8,0.95)(20.8,0)};
      \node[above, orange!80!black, font=\small] at (axis cs:20.55,0.97) {écho B};
    \end{axis}
  \end{tikzpicture}
  \end{center}

  \vspace{4pt}
  \textbf{Données lues sur l'oscillogramme :}
  \begin{itemize}[noitemsep]
    \item \textbf{Écho B} (fond de la pièce) : $\Delta t_B = 20{,}3~\mu\text{s}$
    \item \textbf{Écho A} (anomalie interne) : $\Delta t_A = 8{,}1~\mu\text{s}$
  \end{itemize}

  \vspace{4pt}
  \textit{Rappel : $1~\mu\text{s} = 10^{-6}$~s \hspace{1.5cm} Célérité du son dans l'acier : $v = 5\,900$~m/s}

\end{tcolorbox}

\question{En utilisant l'écho B et la formule $d = \dfrac{v \times \Delta t_B}{2}$, calculer l'épaisseur de la pièce mesurée. Comparer au résultat attendu ($e = 60$~mm).}

\cadrecalcul[2.3cm]

\question{En utilisant l'écho A, calculer la profondeur à laquelle se trouve l'anomalie interne dans la pièce.}

\cadrecalcul[2.3cm]

\question{\textbf{Synthèse :} Le technicien décide de rejeter cette pièce. Justifier cette décision à partir des résultats obtenus aux questions précédentes.}
\reponse{3}

\vspace{8pt}
\begin{center}
  \color{bluedark}\rule{0.5\linewidth}{0.6pt}
\end{center}

\end{document}
